\section{Local star compensation and global trace}
\label{sec:local-star}

Operator moments permit cancellation between closed walks starting at
different rows.  The star certificate of TPC--68 instead tests each row
locally.  We record the exact relation between the two viewpoints.

\subsection{Rowwise compensation identity}

For an active row \(m\), let
\[
 r_m=(R_{mn})_{n\ne m},
 \qquad
 e_m=(E_{mn})_{n\ne m}
\]
be its residual and erasure stars.  Define
\begin{align}
 s_m(R)&:=\|r_m\|_2,
 &
 s_m(E)&:=\|e_m\|_2,
 &
 s_m(H)&:=\|r_m+e_m\|_2,
 \label{eq:local-stars}\\
 c_m&:=2\operatorname{Re}
 \sum_{n\ne m}R_{mn}\overline{E_{mn}}.
 \label{eq:local-compensation}
\end{align}

\begin{proposition}[Exact local compensation]
\label{prop:local-compensation}
For every active row,
\begin{equation}
 \boxed{
 s_m(H)^2=s_m(R)^2+s_m(E)^2+c_m.}
 \label{eq:local-compensation-identity}
\end{equation}
Moreover,
\begin{equation}
 \bigl(s_m(R)-s_m(E)\bigr)^2
 \le s_m(H)^2
 \le
 \bigl(s_m(R)+s_m(E)\bigr)^2.
 \label{eq:local-star-triangle}
\end{equation}
The global mixed trace is the sum of local compensations:
\begin{equation}
 \boxed{
 2\tr(RE)=\sum_m c_m.}
 \label{eq:global-local-compensation}
\end{equation}
\end{proposition}

\begin{proof}
Expand \(\|r_m+e_m\|_2^2\) and apply Cauchy--Schwarz to its cross term.
Since \(R,E\) are Hermitian,
\[
 \tr(RE)
 =
 \sum_{m,n}R_{mn}E_{nm}
 =
 \sum_m\operatorname{Re}
 \sum_{n\ne m}R_{mn}\overline{E_{mn}}.
\]
This proves the final identity.
\end{proof}

Thus a small \(s_m(H)\) in a row where both component stars are large
requires almost extremal negative alignment:
\[
 c_m\approx-2s_m(R)s_m(E).
\]
A cancellation statement for the scalar sum
\(\sum_mc_m\) does not imply this rowwise alignment.

\subsection{What is needed for a star route}

Let
\[
 s_*(H):=\max_m s_m(H)
\]
and let \(\Delta_H\) be the maximum degree of the actual nonzero collision
graph.  TPC--68 proves
\begin{equation}
 s_*(H)\le\|H\|\le\sqrt{\Delta_H}\,s_*(H).
 \label{eq:star-sandwich}
\end{equation}
Consequently, a target \(s_*(H)\le\eta\) requires the uniform inequalities
\begin{equation}
 c_m
 \le
 \eta^2-s_m(R)^2-s_m(E)^2
 \qquad\text{for every active }m.
 \label{eq:uniform-local-compensation-target}
\end{equation}
An average over rows, a density-one set of rows, or a cancellation between
positive \(c_m\)'s and negative \(c_m\)'s does not prove
\eqref{eq:uniform-local-compensation-target}.

The rectangle formula refines the local cross term:
\begin{equation}
 \frac{c_m}{2}
 =
 \operatorname{Re}
 \sum_{\substack{i\ne j\\n\ne m}}
 \one_{\{s_i(m)\ne s_i(n)\}}
 \one_{\{s_j(m)=s_j(n)\}}
 \mathcal H_{i,j}(m,n).
 \label{eq:local-cross-key-rectangles}
\end{equation}
Thus a local star proof requires a cross-key rectangle estimate uniformly
in the base row \(m\), not merely the complete global sum in
\eqref{eq:full-cross-key-rectangle}.

\subsection{Why a normalized global trace is insufficient}

For a zero-diagonal Hermitian matrix,
\begin{equation}
 \tr H^2=\sum_m s_m(H)^2.
 \label{eq:trace-two-star-sum}
\end{equation}
The unnormalized trace is a valid upper certificate:
\[
 \|H\|^2\le\tr H^2.
\]
However, in a row space of growing dimension \(r\), an estimate for the
average \(r^{-1}\tr H^2\) does not bound \(\|H\|\).  One nonzero two-row
block can have norm one while \(r^{-1}\tr H^2\to0\) after isolated rows are
adjoined.

Similarly, \(\tr(RE)=0\) says only that the local compensations sum to
zero.  It permits one component with coherent reinforcement and another
with exact cancellation.  A countermodel is given in
\cref{ex:global-cross-trace-zero}.

For higher even moments the full unnormalized quantity
\[
 \bigl(\tr H^{2q}\bigr)^{1/(2q)}
\]
does bound \(\|H\|\).  But a fixed normalized moment or one isolated sector
need not have the correct scale.  If \(q=q_X\) grows, all constants in the
literal decorated-walk estimate must be uniform in both \(X\) and \(q_X\).
